\section{数値解析：誤差を発生箇所で分ける}
\label{theory:numerics:errors}
物理モデルを計算する際には、材料や境界条件を単純化するモデル誤差、連続な場を格子や多項式で表す離散化誤差、方程式を途中で解くのを止める反復誤差、有限の桁で演算する丸め誤差がある。実験と計算の差をすべてメッシュの粗さへ帰してはいけない。同じ評価量を固定し、どの操作で結果が変わるかを切り分ける。

例えば形状を変えず格子幅を半分にして応力が変わるなら離散化の影響が疑われる。格子を固定して停止閾値を厳しくすると値が変わるなら反復が不足している。倍精度へ変える操作は仮数の精度を増すが、モデルの誤りや切り捨てた高次項を復元しない。反復回数が多いから閾値を緩めるという操作は、速度と引換えに許容誤差を増す操作である。

\subsection{有限精度：桁落ちと情報落ち}
\label{theory:numerics:floating}
浮動小数点数は指数と有限桁の仮数で値を表す。大きな数へ小さい数を足すと、指数を揃えた小さい側の桁が保持範囲外へ落ちることがある。近い二数の差では共通の大きい部分が消え、各入力に既に含まれた誤差が差の値に対して大きくなる。これは異なる現象である。正確に表された二数の減算が常に新しい大誤差を作るわけではない\parencite{goldberg1991}。

実数の結合則は浮動小数点演算の同じ結果を保証しない。通常の倍精度で$10^{16}+1$の1が失われる場合、$(10^{16}+1)-10^{16}$と$(10^{16}-10^{16})+1$は異なる。桁落ちを避ける式変形の例は
\[
 \sqrt{x+1}-\sqrt{x}=\frac{1}{\sqrt{x+1}+\sqrt{x}}\quad(x\ge0)
\]
である。大きい$x$で左辺を近い数の差として直接計算するより、右辺は差による増幅を避けられる。ただしオーバーフローや入力精度など別の限界は残る。

\subsection{差分式：整合性と精度を別に調べる}
\label{theory:numerics:diff}
前節\ref{theory:analysis:difference}のTaylor展開から、同じ一階導関数を近似しても前進・後退差分の誤差は$O(h)$、中心差分は$O(h^2)$となる。二次精度の後退差分も
\[
 \frac{3f_i-4f_{i-1}+f_{i-2}}{2h}=f'_i+O(h^2)
\]
と導ける。定数項の係数和は$3-4+1=0$、一次項は$2hf'_i$、二次項は消える。補間多項式の微分から差分と余項を得る整理もある\parencite{nist_differentiation}。

式が目的の導関数に近づくことが整合性であり、誤差が減る次数が精度である。$(f_{i+1}-2f_i+f_{i-1})/h$は$hf_i''+O(h^3)$であって、一階微分に整合しない。分母を$h^2$に直しても二階微分である。検算には$f=1,x,x^2$を代入するとよい。格子を過度に細かくすると近い値の差の丸め誤差が増し、打切り誤差の低下だけから最終精度を断定できない。

\subsection{補間：与えられた点と元の関数を区別する}
\label{theory:numerics:interpolation}
相異なる$n+1$点を通る次数$n$以下の多項式は一意である。二つあれば差の多項式は$n+1$個の根を持つが次数は$n$以下であるため、差は零となる。Lagrange形式では
\[
 p(x)=\sum_{j=0}^{n}f(x_j)\prod_{k\ne j}\frac{x-x_k}{x_j-x_k}.
\]
各基底は自分の節点で1、他の節点で0となる\parencite{nist_interpolation}。三点なら$p=ax^2+bx+c$へ代入して解く方法も使える。補間は与えられた点を正確に通すが、元の関数全体を正確に復元するとは限らない。区間外の外挿では、補間値の一致だけで信頼性を判断しない。

\subsection{積分領域：三角形の境界を先に書く}
\label{theory:numerics:integration}
頂点$(0,0),(a,0),(0,a)$の三角形は$0\le x\le a,\ 0\le y\le a-x$である。長方形$[0,a]^2$へ置き換えると領域を増やす。一次関数の例では
\[
 \int_0^a\int_0^{a-x}(x+y)\,dy\,dx
 =\int_0^a\left[x(a-x)+\frac{(a-x)^2}{2}\right]dx
 =\frac{a^3}{3}.
\]
独立の幾何検算として、一次関数の三角形内の平均値は重心での値$2a/3$に等しく、面積$a^2/2$との積も$a^3/3$となる。単位も長さの三乗で確認できる。

\subsection{求積：重み和と奇偶の検算}
\label{theory:numerics:quadrature}
前節\ref{theory:analysis:quadrature}のモーメント条件を使うと、$[-1,1]$の一次関数$ax+b$に正確な公式は重み和2かつ$\sum_jw_jx_j=0$を満たす。端点の台形則、中央点則、三点のSimpson則はいずれも満たす。台形則の一般関数に対する余項には二階導関数が現れる\parencite{nist_quadrature}。

例えば端点重みを$1/4$ずつ、中央重みを1とすると重み和は$3/2$であり、定数$b$を$3b/2$へ誤る。$b=0$なら偶然ゼロへ一致するが、全ての一次関数に正確な一般式ではない。特定の値で一致することと、ある関数族に正確な公式であることを分ける。

\subsection{Newton法：更新の符号と終了条件}
\label{theory:numerics:newton}
現在値$x_n$で接線を引き、その零点を次の値とする。$f(x_n)+f'(x_n)(x_{n+1}-x_n)=0$から
\[
 x_{n+1}=x_n-\frac{f(x_n)}{f'(x_n)}.
\]
分母は導関数であり積ではない。単根近傍の条件が満たされると二次収束するが、初期値を問わず収束する保証はなく、$f'(x_n)=0$ではこの更新を使えない\parencite{nist_newton}。

フローチャートが$x\leftarrow x-\Delta x$なら$\Delta x=f/f'$、$x\leftarrow x+\Delta x$なら$\Delta x=-f/f'$である。小さいステップで終了する分岐には$|\Delta x|<\varepsilon$を合わせる。実際の計算ではステップが小さくても残差が十分小さいとは限らないため、$|f(x)|$、相対ステップ、最大反復数、導関数の小ささを別に確認する。試験の簡略図に合わせることと、実務の停止条件が十分であることを分ける。
